% !TEX program = lualatex
\documentclass[lualatex,book,openany,paper=a4paper,jafontsize=11.5pt,jafontscale=1.05]{jlreq}

\usepackage{luatexja-fontspec}
\usepackage{hack-fonts}
\usepackage{hack-cover}
\usepackage[hidelinks]{hyperref}
\usepackage{graphicx}
\usepackage{longtable}
\usepackage{array}
\usepackage{booktabs}
\usepackage{listings}
\usepackage{xcolor}
\usepackage{enumitem}
\usepackage{amsmath}
\usepackage{url}

\setlength{\arrayrulewidth}{0.6pt}
\renewcommand{\arraystretch}{1.18}
\setlength{\emergencystretch}{4em}
\setlist[itemize]{leftmargin=2em}
\setlist[enumerate]{leftmargin=2em}

\lstdefinestyle{textdiagram}{
  basicstyle=\ttfamily\small,
  numbers=none,
  breaklines=true,
  frame=single,
  tabsize=2,
  showstringspaces=false
}

\lstdefinestyle{pseudo}{
  basicstyle=\ttfamily\small,
  numbers=left,
  numberstyle=\tiny,
  breaklines=true,
  frame=single,
  tabsize=2,
  showstringspaces=false
}

\title{ハッカソン}
\subtitle{Processing 音源サブシステム\\基本設計・詳細設計}
\team{チーム23}
\author{
  25G1021 : 梅澤 颯太
}
\date{2026年5月12日}
\version{2.0}

\begin{document}
\hypersetup{pageanchor=false}
\maketitle

\setcounter{tocdepth}{2}
\tableofcontents \thispagestyle{empty}

\clearpage
\hypersetup{pageanchor=true}
\chapter{担当範囲}\setcounter{page}{1}


Processing は，指揮棒の加速度を直接処理しない．加速度から拍・テンポ・強弱を推定する処理は
Arduino 側 node\_01 と node\_02--05 の担当である．Processing 側は，楽器ノードが USB Serial で送る
NOTE パケットを受け取り，\texttt{noteNumber}，\texttt{velocity}，\texttt{durationMs} に従って
音を鳴らす．

\begin{longtable}{p{0.16\linewidth}p{0.32\linewidth}p{0.40\linewidth}}
\caption{担当範囲と本書での対応}\label{tab:wbs}\\
\toprule
WBS & タスク & 本書での設計箇所 \\
\midrule
\endfirsthead
\toprule
WBS & タスク & 本書での設計箇所 \\
\midrule
\endhead
117 & 音色合成 基本設計 & 第2章の機能一覧，第3章のパート対応，第4章の音色設計 \\
124 & NOTE を受け取って鳴らす仕組み & 第4章の NOTE パケット設計，第5章の受信フロー \\
125 & 4 種類の楽器パートを鳴らす管理 & 第3章の partId 対応，第4章の PartManager 設計 \\
126 & 金管音色の具体実装 & 第4章の BrassVoice，倍音比率，ADSR 初期値 \\
254 & UI / モード表示 & 第3章の操作インターフェース，第4章の UiView 設計 \\
\bottomrule
\end{longtable}

\chapter{システムの基本設計}

\section{機能一覧}
表\ref{tab:fbs}に機能一覧を示す．

\begin{longtable}{p{0.28\linewidth}p{0.62\linewidth}}
\caption{Processing 音源サブシステムの機能一覧}\label{tab:fbs}\\
\toprule
機能 & 役割 \\
\midrule
\endfirsthead
\toprule
機能 & 役割 \\
\midrule
\endhead
Serial 接続管理 &
Processing 起動時に Serial ポート一覧を表示し，画面上のポートをクリックして接続する．切断時は再接続待ちへ戻す． \\
NOTE フレーム同期・復号 &
USB Serial のバイト列から \texttt{magic=0x4F52} を探し，20 byte 固定長 NOTE に復号する． \\
パート設定管理 &
期待 \texttt{partId} を \texttt{0x02--0x05} から選び，誤接続時は発音せず警告とログを残す．検証用に ALL モードも持つ． \\
ボイス管理 &
NoteOn，NoteOff，\texttt{durationMs} 経過による自動消音，同時発音数上限，古い Voice の解放を管理する． \\
音色生成 &
\texttt{0x02--0x04} は金管音色，\texttt{0x05} はリズム音色として合成し，Minim の AudioOutput へ出力する． \\
UI 表示・操作 &
接続状態，期待 partId，受信数，破棄数，発音中 Voice 数，波形，警告を表示し，テスト発音やミュートを操作できるようにする． \\
ログ・検証支援 &
受信時刻，NOTE 内容，破棄理由，発音開始を CSV に記録し，結合試験で欠落や誤接続を確認できるようにする． \\
\bottomrule
\end{longtable}

\section{システム構成図}
本番構成は，楽器ノード1台に対して Mac 1台を USB Serial で直結する構成である．
Processing は自分に接続された楽器ノードから NOTE を受け，音声出力を行う．

\begin{lstlisting}[style=textdiagram, caption={Processing 音源サブシステムを含む全体構成}, label={lst:overview}]
User conducts
    |
    v
node_01: conductor, CTRL/BEAT sender
    |
    | UDP multicast CTRL/BEAT
    v
node_02 -- USB Serial NOTE --> Mac 1 Processing final --> Brass 1
node_03 -- USB Serial NOTE --> Mac 2 Processing final --> Brass 2
node_04 -- USB Serial NOTE --> Mac 3 Processing final --> Brass 3
node_05 -- USB Serial NOTE --> Mac 4 Processing final --> Rhythm
\end{lstlisting}

\section{パート対応}
Processing は表\ref{tab:part-map}の対応で \texttt{partId} を解釈する．

\begin{longtable}{p{0.16\linewidth}p{0.14\linewidth}p{0.28\linewidth}p{0.32\linewidth}}
\caption{partId と音色・役割の対応}\label{tab:part-map}\\
\toprule
ノード & partId & 楽譜上の役割 & Processing 側音色 \\
\midrule
\endfirsthead
\toprule
ノード & partId & 楽譜上の役割 & Processing 側音色 \\
\midrule
\endhead
node\_02 & 0x02 & 主旋律 & Brass 1．基準となる金管音色． \\
node\_03 & 0x03 & 輪唱パート1 & Brass 2．金管音色を共有し，音量差で区別する． \\
node\_04 & 0x04 & 輪唱パート2 & Brass 3．金管音色を共有し，必要に応じて調整する． \\
node\_05 & 0x05 & ベース・リズム伴奏 & RhythmVoice．短い減衰を持つリズム音として鳴らす． \\
\bottomrule
\end{longtable}

\chapter{操作インターフェースと状態遷移}

\section{操作インターフェース}
Processing の利用者は，観客ではなく，デモ準備者または結合試験担当者である．
したがって UI は演奏表現ではなく，接続確認，異常検出，安全停止を目的とする．

\begin{longtable}{p{0.22\linewidth}p{0.26\linewidth}p{0.40\linewidth}}
\caption{Processing 側操作インターフェース}\label{tab:ui}\\
\toprule
操作・表示 & 対象 & システム挙動 \\
\midrule
\endfirsthead
\toprule
操作・表示 & 対象 & システム挙動 \\
\midrule
\endhead
起動 & Mac 担当者 & Serial ポート一覧を取得し，画面左に接続候補として表示する． \\
ポートクリック & Mac 担当者 & 選択した Serial ポートを開き，成功したら Ready へ遷移する． \\
\texttt{1}--\texttt{4} キー & Mac 担当者 & 期待 partId を 0x02--0x05 に切り替える．本番では Mac ごとに固定する． \\
\texttt{a} キー & 検証担当者 & ALL モードへ切り替え，0x02--0x05 の全パートを受け付ける． \\
\texttt{r} キー & Mac 担当者 & Serial を閉じ，ポート一覧を再列挙する． \\
\texttt{d} キー & Mac 担当者 & Serial を切断し，PortSelect へ戻す． \\
\texttt{m} キー & Mac 担当者 & ミュートを切り替える．ミュート時は既存 Voice を Release する． \\
\texttt{t} キー & Mac 担当者 & 選択中 partId のテスト音を鳴らし，スピーカと音量を確認する． \\
\texttt{g} キー & 検証担当者 & 疑似 NOTE フレームを内部注入し，Arduino なしで受信処理を確認する． \\
画面表示 & 全員 & 接続状態，期待 partId，受信数，破棄数，欠落数，Voice 数，波形，警告を表示する． \\
\bottomrule
\end{longtable}

\section{状態遷移}
\begin{lstlisting}[style=textdiagram, caption={Processing 音源サブシステムの状態遷移}, label={lst:state}]
Boot
  -> PortSelect  : setup completed and serial ports listed
PortSelect
  -> Ready       : user clicks a serial port and open succeeds
  -> Error       : serial open fails
Ready
  -> Playing     : first valid NOTE received
  -> PortSelect  : user presses r or d
Playing
  -> Ready       : no NOTE for 1000 ms, serial still open
  -> PortSelect  : user presses r or d
  -> Error       : serial exception
Muted
  -> Playing     : user presses m again
  -> PortSelect  : user presses r or d
Error
  -> PortSelect  : user presses r
\end{lstlisting}

\begin{longtable}{p{0.18\linewidth}p{0.30\linewidth}p{0.40\linewidth}}
\caption{状態定義}\label{tab:states}\\
\toprule
状態 & 意味 & 主な処理 \\
\midrule
\endfirsthead
\toprule
状態 & 意味 & 主な処理 \\
\midrule
\endhead
Boot & 起動直後 & 画面，Minim，AudioOutput，Logger，PartManager，SerialFrameReader を初期化する． \\
PortSelect & Serial ポート選択待ち & ポート一覧を表示し，クリック入力を待つ． \\
Ready & 接続済み・NOTE 待ち & Serial バッファを空にし，NOTE の magic 同期を待つ． \\
Playing & 正常受信・発音中 & NOTE 復号，partId 検査，Voice 生成，ログ出力，画面更新を行う． \\
Muted & 受信継続・発音停止 & 新規 Voice は作らず，既存 Voice を Release する． \\
Error & 接続異常 & Serial を閉じ，未処理キューと受信バッファをクリアし，警告を表示する． \\
\bottomrule
\end{longtable}

\chapter{詳細設計}

\section{ハードウェア設計}
Processing 側では回路を追加しないが，接続される構成要素と故障時動作を明確化する．

\begin{longtable}{p{0.28\linewidth}p{0.12\linewidth}p{0.50\linewidth}}
\caption{Processing 側関連部品}\label{tab:parts}\\
\toprule
要素 & 数量 & 役割 \\
\midrule
\endfirsthead
\toprule
要素 & 数量 & 役割 \\
\midrule
\endhead
Arduino UNO R4 WiFi & 4 & node\_02--05．楽譜進行後，NOTE パケットを USB Serial で送信する． \\
Mac & 4 & Processing 実行機．本番は1台につき1つの Serial ポートを開く． \\
USB Type-C ケーブル & 4 & Arduino と Mac の給電および Serial 通信を兼ねる． \\
スピーカまたは内蔵スピーカ & 4 & Processing の AudioOutput を再生する． \\
Processing 4 + Minim & 4環境 & NOTE を受けて音を合成する実行環境． \\
\bottomrule
\end{longtable}

\begin{longtable}{p{0.24\linewidth}p{0.64\linewidth}}
\caption{Processing 側ハードウェア接続設計}\label{tab:hw-design}\\
\toprule
項目 & 設計 \\
\midrule
\endfirsthead
\toprule
項目 & 設計 \\
\midrule
\endhead
接続方式 & node\_02--05 と Mac 1--4 を USB Type-C で 1 対 1 接続する． \\
通信方式 & USB CDC Serial．Arduino 側は \texttt{Serial.write()} で NOTE を20 byte固定長バイナリとして送る． \\
ボーレート & 115200 bps．Arduino 側の NoteSender 設定と統一する． \\
フレーミング & 全 NOTE は先頭2 byte の \texttt{magic=0x4F52} を同期マークとする．Serial 上では \texttt{0x52 0x4F} の順に現れる． \\
音声出力 & 各 Mac の既定オーディオ出力へ送る．本番前に OS 音量と Processing 側ゲインを確認する． \\
誤接続対策 & Mac ごとに期待 \texttt{partId} を設定し，異なる \texttt{partId} は発音せず警告とログを残す． \\
切断時動作 & Serial 切断または例外時はポートを閉じ，受信キューとバッファをクリアし，全 Voice を停止する． \\
\bottomrule
\end{longtable}

\section{ソフトウェア設計}


\subsection{NOTE パケット設計}
Processing は表\ref{tab:note}の20 byteを固定長 NOTE として読む．
Processing 側で独自形式を作らず，塩澤設計の \texttt{OrcProtocol} に合わせる．

\begin{longtable}{p{0.10\linewidth}p{0.18\linewidth}p{0.16\linewidth}p{0.44\linewidth}}
\caption{NOTE パケットの受信側解釈}\label{tab:note}\\
\toprule
offset & field & type & Processing 側の扱い \\
\midrule
\endfirsthead
\toprule
offset & field & type & Processing 側の扱い \\
\midrule
\endhead
0 & magic & uint16 & 0x4F52 以外なら同期探索を続ける． \\
2 & version & uint8 & 0x01 以外なら破棄し，version error としてログする． \\
3 & type & uint8 & 3 のときだけ NOTE として処理する．CTRL/BEAT は発音しない． \\
4 & seq & uint32 & \texttt{partId} ごとに前回値を持ち，重複・欠落を検出する． \\
8 & timestampMs & uint32 & Arduino 側の送信時刻．ログへ保存し，遅延評価の参考にする． \\
12 & partId & uint8 & 0x02--0x05．期待 partId または ALL と一致する場合だけ発音する． \\
13 & noteNumber & uint8 & MIDI ノート番号として周波数へ変換する．0 は停止相当として扱う． \\
14 & velocity & uint8 & 0--127 を振幅へ変換する．\texttt{MASTER\_GAIN} を掛けて過大音量を避ける． \\
15 & gate & uint8 & 1=NoteOn，0=NoteOff．その他は破棄する． \\
16 & durationMs & uint16 & NoteOn 時の自動停止予定時刻を決める．0 の場合は最小長を使う． \\
18 & reserved & uint8[2] & 現時点では未使用． \\
\bottomrule
\end{longtable}

\subsection{オブジェクト設計}
\begin{longtable}{p{0.22\linewidth}p{0.30\linewidth}p{0.36\linewidth}}
\caption{Processing 側オブジェクト設計}\label{tab:classes}\\
\toprule
クラス・構造体 & 主なフィールド・メソッド & 役割 \\
\midrule
\endfirsthead
\toprule
クラス・構造体 & 主なフィールド・メソッド & 役割 \\
\midrule
\endhead
\texttt{NotePacket} &
NOTE の各フィールド &
20 byte フレームから復号した発音指令を表す． \\
\texttt{SerialFrameReader} &
\texttt{pushByte()}, \texttt{pollPacket()}, \texttt{clear()} &
Serial 入力を蓄積し，\texttt{0x52 0x4F} を探して20 byteを切り出す． \\
\texttt{PacketQueue} &
\texttt{offer()}, \texttt{poll()} &
実装では \texttt{ConcurrentLinkedQueue<NotePacket>} を用いる．Serial スレッドと draw スレッドを分離する受信キュー． \\
\texttt{PartManager} &
\texttt{handleNote()}, \texttt{noteOn()}, \texttt{noteOff()}, \texttt{update()} &
Voice の生成，同時発音上限，duration 経過による自動停止を管理する． \\
\texttt{Voice} &
\texttt{start()}, \texttt{update()}, \texttt{release()}, \texttt{stop()} &
1つの発音単位を表す基底クラス．partId，noteNumber，amp，開始時刻，停止予定時刻を持つ． \\
\texttt{BrassVoice} &
\texttt{attackMs=20}, \texttt{decayMs=40}, \texttt{sustain=0.75}, \texttt{releaseMs=60} &
金管1--3の音色を生成する．倍音テーブルから作った波形を使う． \\
\texttt{RhythmVoice} &
\texttt{rhythmDecay()}, \texttt{rhythmFreq()}, \texttt{rhythmWave()} &
node\_05 のリズム音色を生成する．noteNumber により Kick/Snare/Hi-hat 風に分ける． \\
\texttt{UiView} &
\texttt{draw()}, \texttt{drawPortButton()}, \texttt{portIndexAt()} &
状態表示，クリック式ポート選択，波形モニタ，操作説明を描画する． \\
\texttt{Logger} &
\texttt{logPacket()}, \texttt{logEvent()}, \texttt{close()} &
結合試験用 CSV を保存する． \\
\bottomrule
\end{longtable}

\section{音色設計}
\subsection{金管音色}
金管音色は，基音に複数の倍音を重ねた波形を \texttt{WavetableGenerator.gen10()} で作る．
初期値は表\ref{tab:harmonics}とする．

\begin{table}[htbp]
\centering
\caption{金管音色の初期倍音比率}\label{tab:harmonics}
\begin{tabular}{cccccccc}
\toprule
倍音 & 1 & 2 & 3 & 4 & 5 & 6 & 7 \\
\midrule
振幅比 & 1.00 & 0.70 & 0.50 & 0.30 & 0.20 & 0.10 & 0.05 \\
\bottomrule
\end{tabular}
\end{table}

ADSR は表\ref{tab:adsr}の初期値とする．音符全体を単純な直線減衰にすると金管らしさが落ちるため，
Attack，Decay，Sustain，Release を分ける．

\begin{longtable}{p{0.15\linewidth}p{0.15\linewidth}p{0.31\linewidth}p{0.29\linewidth}}
\caption{金管 ADSR 初期値}\label{tab:adsr}\\
\toprule
項目 & 初期値 & 意味 & 調整方針 \\
\midrule
\endfirsthead
\toprule
項目 & 初期値 & 意味 & 調整方針 \\
\midrule
\endhead
Attack & 20 ms & 息の入りの立ち上がり & 遅い場合は10 msへ短縮する． \\
Decay & 40 ms & 最大音量から持続音量へ移る時間 & 音が硬い場合は延長する． \\
Sustain & 0.75 & 音符中の持続音量比 & velocity に応じてスケールする． \\
Release & 60 ms & 息を止めた後の余韻 & 音の切れが悪い場合は40 msへ短縮する． \\
\bottomrule
\end{longtable}

\subsection{リズム音色}
node\_05 はリズム伴奏であるため，金管音色ではなく短い減衰音として扱う．
現在の実装では，noteNumber に応じて周波数，波形，減衰時間を切り替える．

\begin{longtable}{p{0.16\linewidth}p{0.20\linewidth}p{0.14\linewidth}p{0.38\linewidth}}
\caption{リズム音色の設計}\label{tab:rhythm}\\
\toprule
noteNumber & 音色カテゴリ & 減衰時間 & 合成方針 \\
\midrule
\endfirsthead
\toprule
noteNumber & 音色カテゴリ & 減衰時間 & 合成方針 \\
\midrule
\endhead
36--37 & Kick/Bass & 150 ms & 低い SINE 波でベース感を作る． \\
38--40 & Snare & 95 ms & 中域の SQUARE 波で短い打音にする． \\
41--44 & Hi-hat & 45 ms & 高域の SAW 波で鋭いリズム音にする． \\
その他 & Click & 70 ms & 未定義音でも無音にならないよう短いクリック音にする． \\
\bottomrule
\end{longtable}

\chapter{処理フローの設計}

\section{起動フロー}
\begin{enumerate}
  \item \texttt{setup()} で画面，Minim，AudioOutput，設定値，Logger，SerialFrameReader，PartManager，UiView を初期化する．
  \item \texttt{lastSeqByPart} を \texttt{-1} で初期化し，part別 seq 管理を開始できる状態にする．
  \item \texttt{Serial.list()} でポート一覧を取得し，PortSelect 状態にする．
  \item ユーザがポート一覧をクリックしたら \texttt{connectSerial(index)} を呼ぶ．
  \item 接続成功なら Ready，失敗なら Error とし，画面とログに理由を残す．
\end{enumerate}

\section{NOTE 受信フロー}
\begin{lstlisting}[style=pseudo, caption={NOTE 受信から発音までの処理フロー}, label={lst:receive-flow}]
serialEvent(serial):
  while serial has bytes:
    frameReader.pushByte(serial.read())
    while frameReader.pollPacket() returns packet:
      packetQueue.offer(packet)

draw():
  drainPackets()
  partManager.update()
  ui.draw()

drainPackets():
  while packetQueue has packet:
    handlePacket(packet)
\end{lstlisting}

\section{パケット検査フロー}
\begin{lstlisting}[style=pseudo, caption={handlePacket の検査フロー}, label={lst:packet-flow}]
handlePacket(packet):
  if version != 1: drop("version error")
  if type != NOTE: drop("non NOTE type")
  if gate is not 0 and not 1: drop("gate error")
  if partId is not 0x02..0x05: drop("unknown part")
  if expectedPartId != ALL and partId != expectedPartId:
    drop("wrong part")
  if seq equals lastSeqByPart[partId]:
    drop("duplicate seq")
  if seq skips lastSeqByPart[partId] + 1:
    missingPackets += skipped count
  lastSeqByPart[partId] = seq
  logger.logPacket("note", packet)
  partManager.handleNote(packet)
\end{lstlisting}

\section{ボイス管理フロー}
\begin{enumerate}
  \item \texttt{gate=1} かつ \texttt{noteNumber>0} かつ \texttt{velocity>0} の場合，NoteOn として扱う．
  \item \texttt{gate=0}，\texttt{noteNumber=0}，\texttt{velocity=0} の場合，NoteOff 相当として扱う．
  \item \texttt{partId=0x02--0x04} は \texttt{BrassVoice}，\texttt{partId=0x05} は \texttt{RhythmVoice} を生成する．
  \item \texttt{velocity} を 0.0--1.0 に正規化し，\texttt{MASTER\_GAIN=0.55} を掛けて振幅を決める．
  \item 同一 partId の非 release Voice が \texttt{MAX\_VOICES\_PER\_PART=4} 以上なら，最も古い Voice を release する．
  \item \texttt{durationMs} が小さすぎる場合は \texttt{MIN\_DURATION\_MS=120} を使う．
  \item \texttt{update()} で停止予定時刻を過ぎた Voice を release し，release 完了後に \texttt{unpatch} して配列から削除する．
\end{enumerate}

\section{例外処理}
\begin{longtable}{p{0.24\linewidth}p{0.36\linewidth}p{0.28\linewidth}}
\caption{Processing 側例外処理}\label{tab:exceptions}\\
\toprule
条件 & 処理 & 目的 \\
\midrule
\endfirsthead
\toprule
条件 & 処理 & 目的 \\
\midrule
\endhead
magic 不一致 & 1 byte ずつ破棄して次の \texttt{0x52 0x4F} を探す． & バイトずれから復帰する． \\
version 不一致 & パケットを破棄し，CSV に理由を残す． & 旧仕様・別仕様との混在を防ぐ． \\
type 不一致 & NOTE 以外は破棄する． & CTRL/BEAT を誤って発音しない． \\
gate 不正 & パケットを破棄する． & 不正値による無限発音を防ぐ． \\
partId 不明 & パケットを破棄する． & 未定義パートを誤発音しない． \\
partId 不一致 & 発音せず画面警告とログを残す． & Mac と Arduino の誤接続を発見する． \\
seq 重複 & 発音せずログのみ残す． & 二重発音を防ぐ． \\
seq 欠落 & 発音は継続し，欠落数を加算する． & 聴感破綻を避けつつ検証可能にする． \\
Serial 切断 & ポートを閉じ，受信キューとバッファをクリアする． & 再接続時の古いデータ混入を防ぐ． \\
\bottomrule
\end{longtable}

\chapter{テスト設計}

\section{テスト設計}
\begin{longtable}{p{0.26\linewidth}p{0.42\linewidth}p{0.22\linewidth}}
\caption{Processing 側テスト設計}\label{tab:tests}\\
\toprule
試験項目 & 手順 & 合否基準 \\
\midrule
\endfirsthead
\toprule
試験項目 & 手順 & 合否基準 \\
\midrule
\endhead
ビルド確認 & \texttt{processing-java --build} を実行する． & エラーなく完了する． \\
NOTE 復号単体 & \texttt{g} キーで疑似 NOTE を注入する． & 受信数が増え，テスト音が鳴る． \\
ポート選択 UI & Serial ポート一覧をクリックする． & 接続状態が Ready になり，ポート名が表示される． \\
partId 誤接続 & 期待 partId と違う NOTE を入力する． & 発音せず wrong part と表示・ログ出力する． \\
NoteOn 自動停止 & \texttt{durationMs=500} の NOTE を入力する． & 約500 ms後に release され，鳴りっぱなしにならない． \\
同時発音上限 & 同一 partId で5音以上を連続入力する． & 古い Voice が release され，上限4を維持する． \\
リズム音色 & partId 0x05 で noteNumber 36，38，42 を入力する． & 短いリズム音として鳴る． \\
Serial 再接続 & \texttt{d} または \texttt{r} を押して再接続する． & キューとバッファがクリアされ，再接続できる． \\
結合演奏 & node\_02--05 と Mac 4台を接続し演奏する． & 欠落なく鳴り，誤接続警告が出ない． \\
\bottomrule
\end{longtable}

\chapter{生成AIの利用について}
本設計書の作成にあたり，生成AIであるCodex CLIを利用した．利用目的とその箇所は，レビュー資料の読み取り，
塩澤設計の解説，章立ての整理，文章案の作成である．

ただし，採用する前提条件，NOTE パケット仕様，Mac 4台構成，partId 対応，音色パラメータ，
テスト項目については，妥当性を確認した上で判断した．
生成AIの出力はそのまま採用せず，塩澤・齋藤の仕様および担当範囲と矛盾しないよう修正した．

\end{document}
